\chapter{النهايات والاتصال}\label{ch:g12:limcont}

يمدّد هذا الفصل مفهوم النهاية من المتتاليات إلى دوال المتغير
الحقيقي، ويقدّم \omterm{def:g12:limcont:continuity}{الاتصال}. والنتيجة المركزية هي
مبرهنة القيم الوسطى، التي تحوّل معلومة نوعية (``$f$
\omterm{def:g12:limcont:continuity}{متصلة} وتغيّر إشارتها'') إلى وجود حلول لمعادلات.

\section{نهايات الدوال}

في هذا الفصل كله، $f$ \omterm{def:g11:func:function}{دالة} معرَّفة على \omterm{def:g10:numbers:interval}{فترة} أو على \omterm{def:g10:numbers:interunion}{اتحاد} فترات
$D \subseteq \R$.

\begin{definition}[النهاية عند ما لا نهاية]\label{def:g12:limcont:liminf}
لتكن $f$ معرَّفة على \omterm{def:g10:numbers:interval}{فترة} $\intco{a}{+\infty}$ وليكن $\ell \in \R$. نقول
إن $f$ \emph{تؤول إلى $\ell$ عند $+\infty$}\index{نهاية!دالة}،
ونكتب $\lim\limits_{x\to+\infty} f(x) = \ell$، إذا كانت كل \omterm{def:g10:numbers:interval}{فترة} مفتوحة
تحوي $\ell$ تحوي كل القيم $f(x)$ من أجل $x$ كبير بما يكفي:
أي إنه من أجل كل $\varepsilon > 0$ يوجد $A$ بحيث يكون من أجل كل $x \geq A$،
$\abs{f(x) - \ell} \leq \varepsilon$.

ونقول إن $f$ تؤول إلى $+\infty$ عند $+\infty$ إذا وُجد، من أجل كل $M \in \R$،
عدد $A$ بحيث يكون $f(x) \geq M$ من أجل كل $x \geq A$. وتُعرَّف النهايات عند $-\infty$
والنهايات المساوية $-\infty$ بصورة نظيرة.
\end{definition}

\begin{definition}[النهاية عند نقطة]\label{def:g12:limcont:limpoint}
ليكن $a \in \R$ ولتكن $f$ معرَّفة بجوار $a$ (إلا ربما عند $a$). نقول
إن $f$ تؤول إلى $\ell \in \R$ عند $a$، ونكتب
$\lim\limits_{x \to a} f(x) = \ell$، إذا وُجد من أجل كل $\varepsilon > 0$
عدد $\delta > 0$ بحيث يكون من أجل كل $x \in D$ حيث
$\abs{x - a} \leq \delta$، $\abs{f(x) - \ell} \leq \varepsilon$.

ونقول إن $f$ تؤول إلى $+\infty$ عند $a$ إذا وُجد من أجل كل $M$ عدد $\delta > 0$
بحيث يكون $f(x) \geq M$ كلما كان $x \in D$ و $\abs{x - a} \leq \delta$.
وتقصر النهايات من جهة واحدة ($x \to a^+$، $x \to a^-$) المتغير $x$ على $x > a$ أو
$x < a$.
\end{definition}

\begin{example}
$\lim\limits_{x\to+\infty} \frac{1}{x} = 0$،
$\lim\limits_{x\to 0^+} \frac{1}{x} = +\infty$،
$\lim\limits_{x\to 0^-} \frac{1}{x} = -\infty$. أما \omterm{def:g11:func:function}{الدالة}
$x \mapsto \frac1x$ فليست لها نهاية (من الجهتين) عند $0$.
\end{example}

\begin{definition}[المقاربات]\label{def:g12:limcont:asymptote}
يكون المستقيم $y = \ell$ \emph{مقاربًا أفقيًا}\index{مقارب} \omterm{def:g10:functions:graph}{لمنحنى}
$f$ عند $+\infty$ (أو $-\infty$) إذا كان
$\lim_{x\to+\infty} f(x) = \ell$ (أو عند $-\infty$). ويكون المستقيم $x = a$
\emph{مقاربًا شاقوليًا}\index{مقارب} إذا كان \omterm{def:g11:func:function}{للدالة} $f$ نهاية لانهائية من جهة واحدة عند $a$.
\end{definition}

\begin{omfigure}
\begin{tikzpicture}
\begin{axis}[omaxis,
  width=8.8cm, height=6cm,
  xmin=-2.9, xmax=4.9, ymin=-5.2, ymax=9.2,
  xtick={1}, ytick={2}]
  \draw[omProp, dashed] (axis cs:-2.9,2) -- (axis cs:4.9,2);
  \draw[omProp, dashed] (axis cs:1,-5.2) -- (axis cs:1,9.2);
  \addplot[omDef, thick, domain=-2.8:0.855, samples=60] {2 + 1/(x-1)};
  \addplot[omDef, thick, domain=1.145:4.8, samples=60] {2 + 1/(x-1)};
\end{axis}
\end{tikzpicture}

{\small \omterm{def:g10:functions:graph}{منحنى} $f(x) = 2 + \frac{1}{x-1}$ مع مقاربه الأفقي
$y = 2$ (عند $\pm\infty$) ومقاربه الشاقولي $x = 1$.}
\end{omfigure}

\begin{proposition}[العمليات على النهايات]\label{prop:g12:limcont:operations}
تصح قواعد \cref{prop:g12:seq:operations} (المجاميع والجداءات وخارج القسمة)
حرفيًا من أجل نهايات الدوال، عند نقطة أو عند ما لا نهاية، مع الصور
غير المحددة الأربع نفسها $(+\infty)+(-\infty)$ و $0\times\infty$ و
$\frac{\infty}{\infty}$ و $\frac{0}{0}$.
\end{proposition}

\begin{proof}
البراهين مطابقة لحالة المتتاليات، بتعويض ``من أجل $n \geq N$''
بالعبارة ``من أجل $x \geq A$'' أو ``من أجل $\abs{x-a} \leq \delta$''.
\end{proof}

\begin{theorem}[نهاية دالة مركَّبة]\label{thm:g12:limcont:composition}
لتكن $f, g$ دالتين وليكن $a, b, c$ أعدادًا حقيقية أو $\pm\infty$. إذا كان
$\lim\limits_{x \to a} f(x) = b$ و $\lim\limits_{y \to b} g(y) = c$، فإن
\[
\lim_{x \to a} g\bigl(f(x)\bigr) = c .
\]
وتصح العبارة نفسها \omterm{def:g12:seq:sequence}{بمتتالية} محل $f$: إذا كان
$u_n \to b$ و $\lim_{y\to b} g(y) = c$، فإن $g(u_n) \to c$.
\end{theorem}

\begin{proof}
خذ $b, c$ منتهيين (والحالات الأخرى نظيرة). ليكن $\varepsilon > 0$.
يوجد $\delta > 0$ بحيث يستلزم $\abs{y - b} \leq \delta$ أن
$\abs{g(y) - c} \leq \varepsilon$؛ ويوجد $\delta' > 0$ بحيث يستلزم
$\abs{x - a} \leq \delta'$ أن $\abs{f(x) - b} \leq \delta$. وبربط
الاثنين نجد $\abs{g(f(x)) - c} \leq \varepsilon$ من أجل
$\abs{x - a} \leq \delta'$.
\end{proof}

\begin{theorem}[المقارنة والحصر]\label{thm:g12:limcont:squeeze}
لتكن $f, g, h$ دوالًا وليكن $a \in \R \cup \{\pm\infty\}$.
\begin{enumerate}
  \item إذا كان $f \leq g$ بجوار $a$ و $f(x) \to +\infty$ لما $x \to a$، فإن
        $g(x) \to +\infty$.
  \item إذا كان $f \leq g \leq h$ بجوار $a$ وآلت $f, h$ إلى النهاية المنتهية
        نفسها $\ell$ عند $a$، فإن $g(x) \to \ell$ لما $x \to a$.
\end{enumerate}
\end{theorem}

\begin{proof}
الحجة نفسها كما في حالة المتتاليات (\cref{thm:g12:seq:squeeze}).
\end{proof}

\begin{method}[حساب النهايات عمليًا]\label{met:g12:limcont:limits}
\begin{itemize}
  \item كثيرات الحدود والدوال الناطقة عند $\pm\infty$: أخرج أعلى قوة
        للمتغير $x$ عاملًا مشتركًا؛ فالنهاية هي نهاية نسبة الحدين
        المهيمنين.
  \item الصورة $\frac{0}{0}$ عند نقطة $a$: أخرج $(x-a)$ عاملًا مشتركًا من البسط ومن
        المقام، أو تعرّف على \omterm{def:g11:deriv:derivative}{مشتقة}
        $\lim_{x\to a}\frac{f(x)-f(a)}{x-a} = f'(a)$ (\cref{ch:g12:deriv}).
  \item الجذور التربيعية: اضرب في العبارة المرافقة.
  \item العوامل المتذبذبة ($\cos$ و $\sin$): استعمل مبرهنة الحصر.
\end{itemize}
\end{method}

\section{الاتصال}

\begin{definition}[الاتصال]\label{def:g12:limcont:continuity}
لتكن $f$ معرَّفة على \omterm{def:g10:numbers:interval}{فترة} $I$ وليكن $a \in I$. تكون \omterm{def:g11:func:function}{الدالة} $f$
\emph{متصلة عند $a$}\index{اتصال} إذا كان
$\lim\limits_{x \to a} f(x) = f(a)$. وتكون \emph{متصلة على $I$}\index{دالة متصلة} إذا كانت
متصلة عند كل نقطة من $I$.
\end{definition}

وحدسيًا، يمكن رسم \omterm{def:g10:functions:graph}{منحنى} \omterm{def:g12:limcont:continuity}{دالة متصلة} على \omterm{def:g10:numbers:interval}{فترة}
دون رفع القلم.

\begin{proposition}[اتصال الدوال المألوفة]\label{prop:g12:limcont:usual}
كثيرات الحدود والدوال الناطقة و $x \mapsto \sqrt{x}$ و $x \mapsto \abs{x}$
و $\cos$ و $\sin$ و $\exp$ و $\ln$ \omterm{def:g12:limcont:continuity}{متصلة} على كل \omterm{def:g10:numbers:interval}{فترة} محتواة
في \omterm{def:g11:func:function}{مجموعة تعريفها}. ومجاميع الدوال \omterm{def:g12:limcont:continuity}{المتصلة} وجداءاتها وخارج قسمتها (حيث تُعرَّف) ومركباتها
\omterm{def:g12:limcont:continuity}{متصلة}.
\end{proposition}

\begin{proof}[برهان جزئي]
ينتج الاستقرار بالجمع والضرب والقسمة والتركيب من
المبرهنات الموافقة على النهايات. \omterm{def:g12:limcont:continuity}{واتصال} $x \mapsto x$ والدوال
الثابتة مباشر من التعريف، ومنه فإن كثيرات الحدود والدوال الناطقة
\omterm{def:g12:limcont:continuity}{متصلة}. أما \omterm{def:g12:limcont:continuity}{اتصال} بقية الدوال المألوفة
فمبرهَن عليه في الفصول التي تُبنى فيها، أو مقبول.
\end{proof}

\begin{example}
\emph{\omterm{def:g11:func:function}{دالة} الجزء الصحيح} $x \mapsto \floor{x}$ (وهي أكبر \omterm{def:g10:numbers:sets}{عدد صحيح}
$\leq x$) ليست \omterm{def:g12:limcont:continuity}{متصلة} عند أي \omterm{def:g10:numbers:sets}{عدد صحيح} $n$: فنهايتها اليسرى هناك هي
$n - 1$ بينما $\floor{n} = n$.
\end{example}

\begin{omfigure}
\begin{tikzpicture}
\begin{axis}[omaxis,
  width=7.2cm, height=5.6cm,
  xmin=-1.9, xmax=3.5, ymin=-2.5, ymax=2.7,
  xtick={-1,1,2,3}, ytick={-2,-1,1,2}]
  \draw[omDef, thick] (axis cs:-1.5,-2) -- (axis cs:-1,-2);
  \draw[omDef, thick] (axis cs:-1,-1) -- (axis cs:0,-1);
  \draw[omDef, thick] (axis cs:0,0) -- (axis cs:1,0);
  \draw[omDef, thick] (axis cs:1,1) -- (axis cs:2,1);
  \draw[omDef, thick] (axis cs:2,2) -- (axis cs:3,2);
  \fill[omDef] (axis cs:-1,-1) circle (1.5pt);
  \fill[omDef] (axis cs:0,0) circle (1.5pt);
  \fill[omDef] (axis cs:1,1) circle (1.5pt);
  \fill[omDef] (axis cs:2,2) circle (1.5pt);
  \draw[omDef, fill=white] (axis cs:-1,-2) circle (1.5pt);
  \draw[omDef, fill=white] (axis cs:0,-1) circle (1.5pt);
  \draw[omDef, fill=white] (axis cs:1,0) circle (1.5pt);
  \draw[omDef, fill=white] (axis cs:2,1) circle (1.5pt);
  \draw[omDef, fill=white] (axis cs:3,2) circle (1.5pt);
\end{axis}
\end{tikzpicture}

{\small تقفز \omterm{def:g11:func:function}{دالة} الجزء الصحيح عند كل \omterm{def:g10:numbers:sets}{عدد صحيح}: فالقيمة (النقطة المملوءة)
تختلف عن النهاية اليسرى (النقطة المفرَّغة).}
\end{omfigure}

\begin{proposition}[المتتاليات والاتصال]\label{prop:g12:limcont:seqcont}
إذا كان $u_n \to \ell$ وكانت $f$ \omterm{def:g12:limcont:continuity}{متصلة} عند $\ell$، فإن $f(u_n) \to f(\ell)$.
وخاصة، إذا تقاربت \omterm{def:g12:seq:sequence}{متتالية} معرَّفة بالعلاقة $u_{n+1} = f(u_n)$ إلى
$\ell$ وكانت $f$ \omterm{def:g12:limcont:continuity}{متصلة} عند $\ell$، فإن $f(\ell) = \ell$.
\end{proposition}

\begin{proof}
الادعاء الأول هو \cref{thm:g12:limcont:composition} مطبَّقة على
\omterm{def:g12:seq:sequence}{المتتالية} $(u_n)$. وأما الثاني، فانتقل إلى النهاية في طرفي
$u_{n+1} = f(u_n)$: فالطرف الأيسر يؤول إلى $\ell$، والطرف الأيمن إلى
$f(\ell)$، والنهايات وحيدة.
\end{proof}

\section{مبرهنة القيم الوسطى}

\begin{theorem}[مبرهنة القيم الوسطى]\label{thm:g12:limcont:ivt}
\index{مبرهنة القيم الوسطى}
لتكن $f$ \omterm{def:g12:limcont:continuity}{متصلة} على $\intcc{a}{b}$ وليكن $k$ أي عدد محصور بين
$f(a)$ و $f(b)$. عندئذٍ يوجد $c \in \intcc{a}{b}$ بحيث
$f(c) = k$.
\end{theorem}

\begin{omfigure}
\begin{tikzpicture}
\begin{axis}[omaxis,
  width=8.8cm, height=6cm,
  xmin=-0.4, xmax=4.4, ymin=-0.5, ymax=4.7,
  xtick={0.5,3.19,3.8}, xticklabels={$a$,$c$,$b$},
  ytick={1.26,2.5,3.98}, yticklabels={$f(a)$,$k$,$f(b)$}]
  \addplot[omProp, dashed, domain=0:3.19] {2.5};
  \addplot[omDef, thick, domain=0.5:3.8] {0.08*x^3 - 0.5*x + 1.5};
  \draw[black, dashed, thin] (axis cs:0.5,0) -- (axis cs:0.5,1.26)
    -- (axis cs:0,1.26);
  \draw[black, dashed, thin] (axis cs:3.8,0) -- (axis cs:3.8,3.98)
    -- (axis cs:0,3.98);
  \draw[black, dashed, thin] (axis cs:3.19,2.5) -- (axis cs:3.19,0);
  \fill (axis cs:0.5,1.26) circle (1.5pt);
  \fill (axis cs:3.8,3.98) circle (1.5pt);
  \fill (axis cs:3.19,2.5) circle (1.5pt);
\end{axis}
\end{tikzpicture}

{\small مبرهنة القيم الوسطى: لا بد \omterm{def:g10:functions:graph}{لمنحنى} متصل يصل
$(a, f(a))$ بالنقطة $(b, f(b))$ من أن يقطع كل مستقيم أفقي
$y = k$ يقع بين $f(a)$ و $f(b)$.}
\end{omfigure}

\begin{proof}
بتعويض $f$ \omterm{def:g11:func:function}{بالدالة} $f - k$ (وبالدالة $k - f$ عند اللزوم)، يمكننا أن نفترض
$f(a) \leq 0 \leq f(b)$ وأن نبحث عن صفر \omterm{def:g11:func:function}{للدالة} $f$. ونبني متتاليتين
\emph{بالتنصيف}\index{تنصيف}: ضع $a_0 = a$ و $b_0 = b$؛
وإذا أُعطي $a_n \leq b_n$ مع $f(a_n) \leq 0 \leq f(b_n)$، فليكن
$m = \frac{a_n + b_n}{2}$ ونعرّف
\[
(a_{n+1}, b_{n+1}) =
\begin{cases}
(a_n, m) & \text{إذا كان } f(m) \geq 0,\\
(m, b_n) & \text{إذا كان } f(m) < 0.
\end{cases}
\]
وفي الحالتين يكون $f(a_{n+1}) \leq 0 \leq f(b_{n+1})$، \omterm{def:g12:seq:sequence}{والمتتالية} $(a_n)$
\omterm{def:g12:seq:monotonic}{متزايدة}، و $(b_n)$ \omterm{def:g10:functions:variations}{متناقصة}، و
$b_n - a_n = \frac{b-a}{2^n} \to 0$: فالمتتاليتان متجاورتان (انظر
\cref{exo:g12:seq:10}) وتتقاربان إلى نهاية مشتركة
$c \in \intcc{a}{b}$. وبالاتصال، $f(a_n) \to f(c)$ و
$f(b_n) \to f(c)$؛ وبالانتقال إلى النهاية في $f(a_n) \leq 0$ و
$f(b_n) \geq 0$ نجد $f(c) \leq 0 \leq f(c)$، إذن $f(c) = 0$.
\end{proof}

\begin{remark}
برهان التنصيف خوارزمية: فهو يعطي تقريبات دقيقة كما نشاء
لحل ما، ويُنصّف الخطأ عند كل خطوة. وهكذا يمكن لحاسوب أن
يحل $f(x) = 0$ وهو لا يعرف إلا كيف يحسب قيم $f$.
\end{remark}

\begin{theorem}[مبرهنة التقابل]\label{thm:g12:limcont:bijection}
لتكن $f$ \omterm{def:g12:limcont:continuity}{متصلة} و\emph{\omterm{def:g12:seq:monotonic}{رتيبة} تمامًا} على \omterm{def:g10:numbers:interval}{فترة} $I$. عندئذٍ
يكون \omterm{def:g10:algebra:equation}{للمعادلة} $f(x) = k$، من أجل كل $k$ محصور تمامًا بين نهايتي $f$ عند طرفي $I$،
حل \emph{وحيد} في $I$.
\end{theorem}

\begin{proof}
الوجود: اختر $a, b \in I$ بحيث يقع $k$ بين $f(a)$ و $f(b)$ (وهذا ممكن
لأن قيم $f$ تقترب من نهاياتها عند الطرفين)، ثم طبّق
مبرهنة القيم الوسطى على $\intcc{a}{b}$. والوحدانية: إذا كان $x < y$، أعطت
الرتابة التامة أن $f(x) \neq f(y)$، إذن لا يمكن أن تكون نقطتان مختلفتان
حلَّين معًا.
\end{proof}

\begin{method}[حل $f(x) = k$ بمبرهنة التقابل]\label{met:g12:limcont:bijection}
لتبرهن على أن \omterm{def:g10:algebra:equation}{للمعادلة} $f(x) = k$ حلًا واحدًا بالضبط في \omterm{def:g10:numbers:interval}{فترة}:
\begin{enumerate}
  \item أنشئ جدول تغيرات $f$ (إشارة $f'$، انظر
        \cref{ch:g12:deriv})؛
  \item على كل \omterm{def:g10:numbers:interval}{فترة} تكون $f$ فيها \omterm{def:g12:limcont:continuity}{متصلة} \omterm{def:g12:seq:monotonic}{ورتيبة} تمامًا،
        قارن $k$ بالقيم أو النهايات عند الطرفين؛
  \item طبّق مبرهنة التقابل على كل \omterm{def:g10:numbers:interval}{فترة} كهذه وعُدّ
        الحلول.
\end{enumerate}
ثم تحدّد آلة حاسبة أو التنصيف موضع كل حل بالدقة المطلوبة.
\end{method}

\begin{example}\label{ex:g12:limcont:cubic}
\omterm{def:g10:algebra:equation}{للمعادلة} $x^3 + x = 1$ حل حقيقي وحيد. وفعلًا،
$f(x) = x^3 + x$ \omterm{def:g12:limcont:continuity}{متصلة} \omterm{def:g12:seq:monotonic}{ومتزايدة} تمامًا على $\R$ (بوصفها مجموع
دالتين متزايدتين تمامًا)، و $\lim_{-\infty} f = -\infty$ و
$\lim_{+\infty} f = +\infty$، إذن تنطبق مبرهنة التقابل مع $k = 1$.
وبما أن $f(0.6) = 0.816 < 1$ و $f(0.7) = 1.043 > 1$، فإن الحل يقع في
$\intoo{0.6}{0.7}$.
\end{example}

\section{تمارين}

\begin{exercise}[$\star$]\label{exo:g12:limcont:1}
احسب النهايات التالية:
\[
\lim_{x\to+\infty} \frac{2x^3 - x + 1}{5x^3 + x^2}, \qquad
\lim_{x\to-\infty} \bigl(x^2 + 3x - 1\bigr), \qquad
\lim_{x\to 2^+} \frac{x+1}{x-2}, \qquad
\lim_{x\to+\infty} \frac{\cos x}{x}.
\]
\end{exercise}

\begin{exercise}[$\star$]\label{exo:g12:limcont:2}
لتكن $f(x) = \dfrac{2x^2 - 3x + 1}{x - 1}$، معرَّفة من أجل $x \neq 1$.
\begin{enumerate}
  \item بيّن أن $f(x) = 2x - 1$ من أجل كل $x \neq 1$.
  \item استنتج $\lim\limits_{x \to 1} f(x)$. فهل يمكن تمديد $f$ إلى
        \omterm{def:g12:limcont:continuity}{دالة متصلة} عند $1$؟
\end{enumerate}
\end{exercise}

\begin{exercise}[$\star$]\label{exo:g12:limcont:3}
حدّد مقاربات \omterm{def:g10:functions:graph}{منحنى} \omterm{def:g11:func:function}{الدالة}
$f(x) = \dfrac{3x - 1}{x + 2}$ وارسم \omterm{def:g10:functions:graph}{المنحنى} تقريبيًا.
\end{exercise}

\begin{exercise}[$\star\star$]\label{exo:g12:limcont:4}
احسب
\[
\lim_{x\to+\infty} \bigl(\sqrt{x^2 + x} - x\bigr)
\qquad\text{و}\qquad
\lim_{x\to 0} \frac{\sqrt{1+x} - 1}{x}.
\]
\end{exercise}

\begin{exercise}[$\star\star$]\label{exo:g12:limcont:5}
لتكن $f(x) = x + \sin x$.
\begin{enumerate}
  \item بيّن أن $f(x) \to +\infty$ لما $x \to +\infty$، مع أن $f$ ليست
        \omterm{def:g12:seq:monotonic}{رتيبة} على أي \omterm{def:g10:numbers:interval}{فترة} $\intco{A}{+\infty}$.
  \item بيّن أن $g(x) = \dfrac{x + \sin x}{x - \sin x}$ تؤول إلى $1$ عند
        $+\infty$.
\end{enumerate}
\end{exercise}

\begin{exercise}[$\star\star$]\label{exo:g12:limcont:6}
بيّن أن \omterm{def:g10:algebra:equation}{للمعادلة} $x^3 - 3x + 1 = 0$ ثلاثة حلول حقيقية بالضبط،
وأعطِ من أجل كل منها \omterm{def:g10:numbers:interval}{فترة} طولها $1$ تحويه.
(إرشاد: ادرس تغيرات $f(x) = x^3 - 3x + 1$؛ ومشتقتها هي
$3x^2 - 3$.)
\end{exercise}

\begin{exercise}[$\star\star$]\label{exo:g12:limcont:7}
لتكن $f \colon \intcc{0}{1} \to \intcc{0}{1}$ \omterm{def:g12:limcont:continuity}{متصلة}. بيّن أن \omterm{def:g11:func:function}{للدالة} $f$
\emph{\omterm{pb:g10:functions:1}{نقطة ثابتة}}: أي أنه يوجد $c \in \intcc{0}{1}$ بحيث $f(c) = c$.
(إرشاد: طبّق مبرهنة القيم الوسطى على $g(x) = f(x) - x$.)
\end{exercise}

\begin{exercise}[$\star\star\star$]\label{exo:g12:limcont:8}
تصعد متنزّهة درب جبل، فتنطلق عند 8:00 وتصل عند 20:00. وفي
اليوم التالي تنزل الدرب نفسه، منطلقة عند 8:00 وواصلة عند
20:00. بيّن أن ثمة لحظة من اليوم تكون فيها عند النقطة نفسها بالضبط
من الدرب في اليومين. (إرشاد: أدخل الفرق بين دالتي
الموضع.)
\end{exercise}

\begin{exercise}[$\star\star\star$]\label{exo:g12:limcont:9}
لتكن $f(x) = x^5 + x^3 + x - 1$.
\begin{enumerate}
  \item بيّن أن \omterm{def:g11:func:function}{للدالة} $f$ صفرًا حقيقيًا وحيدًا $\alpha$، وأن
        $\alpha \in \intoo{0}{1}$.
  \item صِف إجراء تنصيف يحسب $\alpha$ بدقة
        $10^{-6}$، وأعطِ عدد التكرارات اللازمة انطلاقًا من
        $\intcc{0}{1}$.
\end{enumerate}
\end{exercise}

\section{مسألة: لا تعبر نهرًا دون أن تبتلّ}

\begin{problem}[{مسألة نهاية الأسبوع --- مبرهنة القيم
الوسطى تفرض النقط الثابتة، والنقط المتقابلة المتساوية الحرارة،
والطاولات الثابتة: ثلاث مبرهنات تبدو مستحيلة من فكرة واحدة}]
\label{pb:g12:limcont:1}
كوّر خريطة لبلد ما وألقها في أي مكان من أرض ذلك البلد: فثمة
نقطة من الخريطة تقع \emph{بالضبط} فوق الموضع الذي
تمثّله. وفي كل لحظة، توجد نقطتان متقابلتان قطريًا
على خط الاستواء لهما \emph{بالضبط} درجة الحرارة نفسها. وأي
طاولة مربعة متأرجحة على أرض متعرجة يمكن دائمًا تثبيتها
\emph{بإدارتها}. ثلاث حيل، ومبرهنة واحدة: فمقدار متصل
لا يمكن أن ينتقل من السالب إلى الموجب دون أن ينعدم
(\cref{thm:g12:limcont:ivt}). وتدرّب هذه المسألة على النهايات، ثم
تبني العجائب الثلاث كلها --- وتعرف بالضبط ما لا تعد به
المبرهنة.

\medskip
\noindent\textbf{الجزء الأول --- التمكّن من النهايات.}
\begin{enumerate}
  \item احسب $\lim_{x \to +\infty} \dfrac{3x^2 - x}{x^2 + 5}$،
        \ $\lim_{x \to +\infty} \dfrac{2x + 1}{x^2 + 1}$، \ و
        $\lim_{x \to +\infty} (x^3 - x^2)$
        (\cref{met:g12:limcont:limits}).
  \item أعطِ مقاربات
        $f(x) = \dfrac{3x - 1}{x + 2}$
        (\cref{def:g12:limcont:asymptote})، بحساب النهايات
        المناسبة.
  \item احسب $\lim_{x \to 3} \dfrac{x^2 - 9}{x - 3}$
        (\omterm{def:g10:algebra:expand}{بالتحليل إلى عوامل}) و
        $\lim_{x \to 0} \dfrac{\sqrt{x + 1} - 1}{x}$
        (بالمرافق).
  \item بمبرهنة الحصر
        (\cref{thm:g12:limcont:squeeze}): احسب
        $\lim_{x \to 0} x \sin\frac1x$ و
        $\lim_{x \to +\infty} \frac{\sin x}{x}$.
  \item بالتركيب والحدود المهيمنة:
        $\lim_{x \to +\infty} \dfrac{\sqrt{4x^2 + x}}{x}$.
\end{enumerate}

\medskip
\noindent\textbf{الجزء الثاني --- جذور مصادَق عليها.}
\begin{enumerate}[resume]
  \item بيّن أن \omterm{def:g10:algebra:equation}{للمعادلة} $x^3 + x = 5$ حلًا واحدًا على الأقل
        في $\intcc{1}{2}$.
  \item بيّن أن الحل وحيد على $\R$
        (\cref{thm:g12:limcont:bijection}).
  \item برهن على العبارة الكلاسيكية: \emph{لكل كثير حدود درجته \omterm{def:g11:func:parity}{فردية}
        \omterm{def:g11:quad:discriminant}{جذر} حقيقي واحد على الأقل}. (اكتب الحجة
        من أجل $p(x) = x^3 + a x^2 + b x + c$: احسب النهايتين
        اللانهائيتين بالحد المهيمن، ثم طبّق مبرهنة القيم الوسطى على
        \omterm{def:g10:numbers:interval}{فترة} كبيرة.)
  \item تحديد موضع \omterm{def:g11:quad:discriminant}{جذر} السؤال 6 \omterm{thm:g12:limcont:ivt}{بالتنصيف}: كم
        تنصيفًا \omterm{def:g10:numbers:interval}{للفترة} $\intcc{1}{2}$ يضمن خطأ دون
        $10^{-6}$ (\cref{exo:g12:limcont:9})؟ وقارن بتضاعف
        الأرقام في طريقة نيوتن
        (\cref{pb:g11:deriv:1}): فكم خطوة \omterm{def:g10:numbers:approx}{تقريبًا} ستحتاج إليها؟
  \item المثال المضاد الذي يحرس المبرهنة:
        تحقق $g(x) = \frac1x$ أن $g(-1) = -1 < 0$ و
        $g(1) = 1 > 0$، ومع ذلك لا تنعدم أبدًا. فأي فرضية من فرضيات
        مبرهنة القيم الوسطى تسقط --- وماذا يعلّمنا المثال عن
        التحقق من الفرضيات؟
\end{enumerate}

\medskip
\noindent\textbf{الجزء الثالث --- ثلاث مبرهنات تبدو
مستحيلة.}
\begin{enumerate}[resume]
  \item مبرهنة \omterm{pb:g10:functions:1}{النقطة الثابتة} على القطعة: إذا كانت
        $f : \intcc{0}{1} \to \intcc{0}{1}$ \omterm{def:g12:limcont:continuity}{متصلة}، فبرهن
        على أن ثمة $c$ تحقق $f(c) = c$. (ادرس
        $g(x) = f(x) - x$ عند الطرفين.)
  \item ترجم السؤال 11 إلى لغة الخريطة المكوّرة: ما الذي يلعب
        دور $f$، ولماذا تقع صورته في
        $\intcc{0}{1}$ (مجازًا: داخل البلد)، وما
        النقطة الواقعة بالضبط فوق موضعها؟ (وثمة
        بُعد واحد مقبول في صمت --- والعبارة الأمينة ذات
        البعدين هي مبرهنة براور، في الكتب
        الجامعية.)
  \item خط الاستواء: لتكن $T$ درجة حرارة \omterm{def:g12:limcont:continuity}{متصلة} على
        دائرة، منظورًا إليها بوصفها \omterm{def:g11:func:function}{دالة} دورية دورها $2\pi$ حيث
        $T(0) = T(2\pi)$. ضع $g(t) = T(t) - T(t + \pi)$.
        احسب $g(0) + g(\pi)$، واستنتج أن $g$ تنعدم
        في مكان ما من $\intcc{0}{\pi}$، ثم صُغ المبرهنة
        المتعلقة بالنقطتين المتقابلتين.
  \item الطاولة المتأرجحة: طاولة مربعة صلبة تمامًا تقف
        على أرض متعرجة \omterm{def:g12:limcont:continuity}{متصلة} (بلا درجات!)؛ فثلاث أرجل
        تلامس دائمًا، وواحدة تحوم على ارتفاع $h(\theta)$
        يتعلق اتصاليًّا بزاوية دوران الطاولة
        $\theta$. فسّر لماذا يفرض الدوران بربع دورة
        أن تغيّر فجوة التحويم إشارتها --- واستنتج أن
        ثمة زاوية $\theta^*$ تثبّت الأرجل الأربع كلها. (وقد نُشرت هذه
        الحجة مبرهنةً حقيقية في سنة 2005؛
        وأصحاب المقاهي عرفوا الحيلة قبل ذلك.)
  \item استعملت متنزّهة \cref{exo:g12:limcont:8} المخطط
        نفسه. صُغ في ثلاثة أسطر القالب المشترك بين الأسئلة 11 و 13
        و 14 والمتنزّهة --- أي \omterm{def:g11:func:function}{دالة} تبني، وماذا تتحقق منه
        عند الطرفين، وماذا تستدعي.
  \item ما \emph{لا} تعطيه مبرهنة القيم الوسطى: فقد جاء الوجود بلا مقابل،
        لكن أي سؤالين يتركهما مفتوحين، وأي أدوات في هذه
        المسألة تجيب عنهما (السؤالان 7 و 9)؟
\end{enumerate}

\medskip
\noindent\textbf{الجزء الرابع --- صور مع مقاربات.}
\begin{enumerate}[resume]
  \item أعد كتابة $f(x) = \dfrac{3x - 1}{x + 2}$ على الصورة
        $3 - \dfrac{7}{x + 2}$ واستعملها لوصف \omterm{def:g10:functions:graph}{المنحنى}
        وصفًا كاملًا: المقاربات، والتغيرات على كل جهة،
        والرسم التقريبي. (وهو \omterm{def:g11:func:reference}{قطع زائد} منسحب --- \omterm{def:g10:functions:graph}{ومنحنى} التكلفة المتوسطة
        في \cref{pb:g10:reffunc:1} كان كذلك.)
  \item \omterm{def:g11:func:function}{الدالة} $f(x) = \dfrac{x^2 - 1}{x - 1}$ غير
        معرَّفة عند $1$. فأي قيمة تُسنَد إلى $f(1)$ تجعلها
        \omterm{def:g12:limcont:continuity}{متصلة} هناك --- ولماذا يسمّي المحلّلون نقطة كهذه
        انقطاعًا \emph{قابلًا للرفع}؟
  \item \omterm{def:g11:func:function}{دالة} الجزء الصحيح $x \mapsto \lfloor x \rfloor$
        (أكبر \omterm{def:g10:numbers:sets}{عدد صحيح} $\leq x$): أين تكون \omterm{def:g12:limcont:continuity}{متصلة}،
        وأين تقفز، ولماذا تسقط خاتمة مبرهنة القيم الوسطى
        من أجلها (أوجد $a < b$ حيث
        $\lfloor a \rfloor < \frac12 < \lfloor b \rfloor$ ومع ذلك
        لا حل \omterm{def:g10:algebra:equation}{للمعادلة} $\lfloor x \rfloor = \frac12$)؟
  \item الخاتمة --- وجها \omterm{def:g12:limcont:continuity}{الاتصال}: الوعد المحلي
        (القيمة $=$ النهاية: فلا انتقال آني) والقوة الشاملة
        (مبرهنة القيم الوسطى: فكل قيمة وسطى تُزار). صنّف
        عجائب الجزء الثالث الأربع تحت الوجه الصحيح، وأعطِ
        شعار عبور النهر اسمه الرياضي الدقيق.
\end{enumerate}
\end{problem}
